% ---------------------------------------------------------------------------
\subsection{Beyond the linear-Gaussian model}\label{sec:repr-nonlinear}
% ---------------------------------------------------------------------------

This subsection and the next address Problem~P7 of \cref{sec:open}. \Cref{thm:overconstraint,prop:fisher-rep} rest on
three simplifications: a linear head, Gaussian features, and a constant Fisher matrix. We remove them in turn. The
first result is elementary, but it organises the rest: a constraint costs no approximation exactly when some
output-optimal student already satisfies it.

\begin{lemma}[Compatibility and the penalty path]\label{lem:penalty}
\Pv\Ck Let $\cL,\cR:\Theta\to[0,\infty)$ attain their minima $\cL^\star$ and $\cR_0$, and for $\gamma>0$ let $\theta_\gamma$ be any
minimiser of $\cL+\gamma\cR$. Define the over-constraint cost $\Delta(\gamma):=\cL(\theta_\gamma)-\cL^\star$ and its
hard-constraint value $\Delta_\infty:=\inf\{\cL(\theta):\cR(\theta)=\cR_0\}-\cL^\star$.
\begin{enumerate}[label=(\roman*),nosep]
\item If $\argmin\cL\cap\argmin\cR\neq\emptyset$, then for every $\gamma>0$ the minimisers of $\cL+\gamma\cR$ are exactly
$\argmin\cL\cap\argmin\cR$. In particular $\Delta\equiv0$.
\item If $\gamma_1<\gamma_2$, then $\cL(\theta_{\gamma_1})\le\cL(\theta_{\gamma_2})$ and $\cR(\theta_{\gamma_1})\ge\cR(\theta_{\gamma_2})$.
\item $0\le\Delta(\gamma)\le\Delta_\infty$ and $\cR(\theta_\gamma)-\cR_0\le\Delta_\infty/\gamma$.
\item If $\Theta$ is a compact metric space and $\cL,\cR$ are lower semicontinuous, then $\Delta(\gamma)\uparrow\Delta_\infty$ as
$\gamma\to\infty$, and $\Delta_\infty=0$ iff $\argmin\cL\cap\argmin\cR\neq\emptyset$.
\end{enumerate}
\end{lemma}
\begin{proof}
(i) If $\theta_0$ lies in both sets, then $\cL(\theta)+\gamma\cR(\theta)\ge\cL^\star+\gamma\cR_0=\cL(\theta_0)+\gamma\cR(\theta_0)$, with
equality iff $\cL(\theta)=\cL^\star$ and $\cR(\theta)=\cR_0$.
(ii) Write $\theta_i=\theta_{\gamma_i}$. Adding the optimality inequalities
$\cL(\theta_i)+\gamma_i\cR(\theta_i)\le\cL(\theta_j)+\gamma_i\cR(\theta_j)$, $\{i,j\}=\{1,2\}$, gives
$(\gamma_2-\gamma_1)(\cR(\theta_2)-\cR(\theta_1))\le0$. The inequality with $i=1$ then gives
$\cL(\theta_1)-\cL(\theta_2)\le\gamma_1(\cR(\theta_2)-\cR(\theta_1))\le0$.
(iii) For any $\theta'$ with $\cR(\theta')=\cR_0$,
$[\cL(\theta_\gamma)-\cL^\star]+\gamma[\cR(\theta_\gamma)-\cR_0]\le\cL(\theta')-\cL^\star$. Both brackets are non-negative. Take the
infimum over $\theta'$.
(iv) Let $\gamma_n\to\infty$. By (iii), $\cR(\theta_{\gamma_n})\to\cR_0$. Along a subsequence $\theta_{\gamma_n}\to\bar\theta$, and
lower semicontinuity gives $\cR(\bar\theta)=\cR_0$ and $\cL(\bar\theta)\le\lim_n\cL(\theta_{\gamma_n})$; the limit exists by (ii).
Hence $\Delta_\infty\le\lim_n\Delta(\gamma_n)$, and (iii) gives equality. The set $\{\cR\le\cR_0\}$ is compact, so $\cL$ attains
$\cL^\star+\Delta_\infty$ on it. If $\Delta_\infty=0$, the point where it does lies in both argmin sets.
\end{proof}
\noindent
\Cref{prop:fisher-rep} is case (i), and \cref{thm:overconstraint}(b) computes $\Delta_\infty$ for $\cR_{\mathrm{LS}}$. So, for any
architecture and any loss, the question ``does this constraint cost approximation?'' is the question ``is it
compatible with some output-optimal student?''.

\paragraph{Nonlinear heads and arbitrary feature laws.}
Keep the linear bottleneck $h_\TS=Bh_\TT$ of \cref{thm:overconstraint}, and make everything else general. The teacher
representation $h_\TT$ is any centred random vector in $\R^{d_\TT}$ with covariance $\Sigma$ and eigenpairs
$(\lambda_i,u_i)$, $\lambda_1\ge\lambda_2\ge\cdots$. The teacher's next-token law is $p_\TT(\cdot\mid s)=\pi(h_\TT(s))$ for an
arbitrary measurable head $\pi:\R^{d_\TT}\to\Delta(\cV)$. The student is $q_\theta(\cdot\mid s)=\kappa(Bh_\TT(s))$ with an arbitrary
measurable head $\kappa:\R^{d_\TS}\to\Delta(\cV)$ (a \emph{universal head}). The loss is the exact forward KL,
$\cL(B,\kappa)=\E\,\KL(\pi(h_\TT)\,\|\,\kappa(Bh_\TT))$. Let $Y$ be a teacher token: $Y\mid h_\TT\sim\pi(h_\TT)$.

\begin{theorem}[Over-constraint for nonlinear teachers]\label{thm:nl-overconstraint}
\Mo\Ck
\begin{enumerate}[label=(\alph*),nosep]
\item $\min_\kappa\cL(B,\kappa)=I(Y;h_\TT\mid Bh_\TT)$, attained by $\kappa(u)=\E[\pi(h_\TT)\mid Bh_\TT=u]$. The output-only optimum
is $\cA_{\mathrm{out}}=\inf_BI(Y;h_\TT\mid Bh_\TT)$. It depends on $B$ only through $\operatorname{range}(B^\top)$.
\item Suppose $\lambda_{d_\TS}>\lambda_{d_\TS+1}$. The minimisers of $\cR_{\mathrm{LS}}(B)=\min_A\E\|ABh_\TT-h_\TT\|^2$ are the $B$ with
$Bh_\TT=MU_{d_\TS}^\top h_\TT$ a.s.\ for an invertible $M$, where $U_{d_\TS}=(u_1,\dots,u_{d_\TS})$. Under full feature matching as a
hard constraint, the best achievable loss is therefore $I(Y;h_\TT\mid U_{d_\TS}^\top h_\TT)$.
\item The over-constraint cost $\Delta_\infty=I(Y;h_\TT\mid U_{d_\TS}^\top h_\TT)-\cA_{\mathrm{out}}$ is $\ge0$, and it is $0$ iff the
principal subspace is output-optimal. It can equal $I(Y;h_\TT)$, which is the loss of a student that ignores its input
and which tends to $\log V$ for a nearly deterministic teacher. Example: $d_\TT=2$, $d_\TS=1$, independent coordinates with
$\Var h_1>\Var h_2$, $h_2$ uniform on $V$ values, and a teacher that reads only $h_2$.
\end{enumerate}
\end{theorem}
\begin{proof}
(a) Let $\bar\pi:=\E[\pi(h_\TT)\mid Bh_\TT]$. For any $Bh_\TT$-measurable $\kappa$,
$\E\KL(\pi\|\kappa)=\E\KL(\pi\|\bar\pi)+\E\KL(\bar\pi\|\kappa)$, because conditioning on $Bh_\TT$ turns
$\E\sum_y\pi_y\log(\bar\pi_y/\kappa_y)$ into $\E\sum_y\bar\pi_y\log(\bar\pi_y/\kappa_y)$. The first term is $I(Y;h_\TT\mid Bh_\TT)\le\log V$,
since $\pi(h_\TT)$ and $\bar\pi$ are the laws of $Y$ given $h_\TT$ and given $Bh_\TT$. If $M$ is invertible, $Bh_\TT$ and $MBh_\TT$
generate the same $\sigma$-algebra.
(b) Only second moments enter. With $P$ the orthogonal projector onto $\operatorname{range}(\Sigma^{1/2}B^\top)$, of dimension
$\le d_\TS$, the normal equations give $\cR_{\mathrm{LS}}(B)=\tr\Sigma-\tr(P\Sigma)$. By Ky Fan's maximum principle,
$\tr(P\Sigma)\le\sum_{i\le d_\TS}\lambda_i$, with equality iff $\operatorname{range}P=\operatorname{span}(u_1,\dots,u_{d_\TS})$; the eigengap forces
$\lambda_{d_\TS}>0$. Then $\Sigma^{1/2}B^\top=U_{d_\TS}N$ with $N$ invertible, and writing $h_\TT=\Sigma^{1/2}\xi$ gives
$Bh_\TT=N^\top\Lambda_{d_\TS}^{-1/2}U_{d_\TS}^\top h_\TT$ with $\Lambda_{d_\TS}=\diag(\lambda_1,\dots,\lambda_{d_\TS})$. Now apply (a).
(c) Non-negativity is (a). In the example the principal direction is $e_1$. Keeping $h_2$ gives $I(Y;h_\TT\mid h_2)=0$, while
$I(Y;h_\TT\mid h_1)=I(Y;h_2)=I(Y;h_\TT)$ by independence. This tends to $H(Y)=\log V$ when $\pi(h_2)$ is nearly a point mass on
distinct tokens.
\end{proof}
\noindent
Neither Gaussianity nor linearity was needed. Full feature matching can destroy \emph{all} the information that the
teacher's output carries about its representation, whatever the head.

\paragraph{Zero cost for nonlinear pairs: the Fisher central subspace.}
Let the teacher's centred logits be $\tilde z_\TT(s)=g(h_\TT(s))$ with $g:\R^{d_\TT}\to\operatorname{range}C$, so that
$\pi=\softmax\circ g$. Let $J(h)=Dg(h)$ and $F(h)=\diag\pi(h)-\pi(h)\pi(h)^\top$. As in \cref{def:fisher-rep}, set
\[
K:=\E\big[J(h_\TT)^\top F(h_\TT)J(h_\TT)\big]=\E\big[\nabla_h\log\pi_Y(h_\TT)\,\nabla_h\log\pi_Y(h_\TT)^\top\big],
\]
the Fisher information that one teacher token carries about the teacher's representation. The identity holds because
$\nabla_h\log\pi_y=J^\top(e_y-\pi)$. $K$ is a Fisher-weighted expected gradient outer product, the object behind average
derivative estimation and active subspaces \citep{samarov1993exploring,constantine2014active}. Let
$\mathcal S_F:=\operatorname{range}K$, let $k_F:=\rank K$, and let $U_F\in\R^{d_\TT\times k_F}$ be an orthonormal basis of $\mathcal S_F$.

\begin{theorem}[The Fisher central subspace]\label{thm:fisher-central}
\Pa{the stated regularity}\Ck Suppose $h_\TT$ has a density that is positive on an open convex set $\Omega$ carrying all of its
mass, $g$ is $C^1$ on $\Omega$, and $\E\|J(h_\TT)\|^2<\infty$.
\begin{enumerate}[label=(\roman*),nosep]
\item For $v\in\R^{d_\TT}$: $v^\top Kv=0$ iff $J(h)v=0$ for every $h\in\Omega$. Consequently $g(h)=g(h')$ whenever $h,h'\in\Omega$ and
$h-h'\perp\mathcal S_F$, that is, $g=\phi\circ U_F^\top$ on $\Omega$ for some function $\phi$. $\mathcal S_F$ is the smallest subspace with
this property. The teacher is a multi-index model whose index space is $\mathcal S_F$, however $J(h)$ varies.
\item (\emph{Zero approximation cost.}) Suppose $k_F\le d_\TS$. Let $t_F(s):=U_F^\top h_\TT(s)$ and
$\cR_{\mathrm F}(\theta):=\min_A\E\|Ah_\TS(s)-t_F(s)\|^2$. Assume \emph{layerwise realisability}: some student has
$h_\TS(s)=(Mt_F(s),0)$ with $M$ invertible, and a head that maps this to $\phi(t_F(s))$. Then for every $\gamma>0$ the minimisers
of $\E\KL(p_\TT\|q_\theta)+\gamma\cR_{\mathrm F}(\theta)$ are exactly the lossless students with $\cR_{\mathrm F}=0$. In the model of
\cref{thm:nl-overconstraint} realisability is automatic. If $\Sigma\succ0$, the target of \cref{def:fisher-rep} built from the
top-$k_F$ eigenvectors of $\Sigma^{1/2}K\Sigma^{1/2}$ is an invertible linear image of $t_F$, so the same conclusion holds for it.
\item (\emph{A pointwise Jacobian target.}) The same conclusion holds for
\[
\cR_J(\theta):=\min_{A\in\R^{d_\TT\times d_\TS}}\E\big\|F(h_\TT)^{1/2}J(h_\TT)\big(Ah_\TS(s)-h_\TT(s)\big)\big\|^2 .
\]
This constraint depends on $s$ through the Jacobian. It is the second-order expansion of the head-transplant loss
$\min_A\E\,\KL(\pi(h_\TT)\,\|\,\pi(Ah_\TS))$, which feeds the student's mapped features through the teacher's head
\citep[cf.][]{chen2022reused}. It needs only Jacobian--vector products through the teacher's upper network, not $K$.
\end{enumerate}
\end{theorem}
\begin{proof}
(i) $v^\top Kv=\E\|F^{1/2}Jv\|^2$. Since $\pi(h)$ has full support, $\ker F(h)=\operatorname{span}(\one)$, while
$J(h)v\in\operatorname{range}C=\one^\perp$. Hence $v^\top Kv=0$ iff $J(h_\TT)v=0$ a.s. The set $\{h\in\Omega:J(h)v\neq0\}$ is open, so it is
null only if it is empty. For $h,h'\in\Omega$ with $h'-h\in\mathcal S_F^\perp=\ker K$, the segment $[h,h']$ lies in $\Omega$, and
$\frac{\mathrm d}{\mathrm dt}g(h+t(h'-h))=J(h+t(h'-h))(h'-h)=0$ along it. Conversely, if $g(h)=g(h')$ whenever $Qh=Qh'$ for an
orthogonal projector $Q$, then $J(h)v=0$ for $v\in\ker Q$, so $\ker Q\subseteq\ker K$ and $\mathcal S_F\subseteq\operatorname{range}Q$.
(ii) The realisable student is lossless and has $\cR_{\mathrm F}=0$, so it lies in $\argmin\cL\cap\argmin\cR_{\mathrm F}$. Apply
\cref{lem:penalty}(i). In the model of \cref{thm:nl-overconstraint}, take $B=(MU_F^\top;0)$ and $\kappa=\softmax\circ\phi\circ M^{-1}$ on the
first $k_F$ coordinates. For the whitened target, the top-$k_F$ eigenvectors $U$ of $\Sigma^{1/2}K\Sigma^{1/2}$ span
$\Sigma^{1/2}\mathcal S_F$, so $\Sigma^{-1/2}U=U_FN$ with $N$ invertible and $U^\top\Sigma^{-1/2}h_\TT=N^\top t_F$.
(iii) With $A=(U_FM^{-1},0)$, $Ah_\TS-h_\TT=-(I-U_FU_F^\top)h_\TT\in\ker K$, which $J(h_\TT)$ annihilates by (i). So $\cR_J=0$ at the
realisable student, and (ii) applies. For the expansion, $\KL(p\,\|\softmax(v+\eta))=\frac12\eta^\top F\eta+O(\|\eta\|^3)$ with
$\eta=g(h+\delta)-g(h)=J\delta+o(\|\delta\|)$.
\end{proof}
\noindent
For language models the state law is discrete, so $h_\TT(s)$ has no density. Then $K$ sees the Jacobian only on the
support, and the conclusion can fail: $g$ may vary between two support points along a direction in which its
derivative vanishes at both. The remedy is to compute $K$ under a smoothed law, for example that of
$h_\TT(s)+\sigma\varepsilon$ with Gaussian $\varepsilon$, whose density is positive on an open convex set containing the support.
Then (i) holds on that set, and (ii) and (iii) hold for the original state law.

So the answer to the first question of P7 is \emph{yes} whenever the teacher's Fisher central subspace fits in the
student. The average $K$ over states does not blur alignment: its range is exactly the set of directions that the
nonlinear head ever uses, and full feature matching can still cost everything (\cref{thm:nl-overconstraint}(c)). When
$k_F>d_\TS$ no student is lossless, and the right target depends on the student's head.

\paragraph{Lossy regime I: linear read-outs, and which average of the Jacobian.}
Use the quadratic loss $\cL=\frac12\E\|F^{1/2}(\tilde z_\TS-\tilde z_\TT)\|^2$ with a constant $F\succeq0$, as in
\cref{thm:overconstraint}; by \cref{prop:temperature}(ii) it is the exact high-temperature limit, with $F=\frac1VC$. The teacher is
$\tilde z_\TT=g(h_\TT)$ with $g$ nonlinear and square integrable, and $h_\TT$ is centred with $\Sigma\succ0$. The student is linear:
$\tilde z_\TS=W_\TS Bh_\TT+b$. Put $\mathsf C:=\Cov(g(h_\TT),h_\TT)$ and $G_\sharp:=F^{1/2}\mathsf C\Sigma^{-1/2}$, with singular values
$\sigma_1\ge\sigma_2\ge\cdots$ and right singular vectors $v_1,v_2,\dots$. Let
$\mathcal N:=\frac12\E\|F^{1/2}(g(h_\TT)-\E g-\mathsf C\Sigma^{-1}h_\TT)\|^2$ be the part of the teacher that no linear student can fit.

\begin{theorem}[Linear read-outs: average the Jacobian, not the Fisher]\label{thm:stein}
\Mo\Ck
\begin{enumerate}[label=(\alph*),nosep]
\item $\min\cL=\mathcal N+\frac12\sum_{i>d_\TS}\sigma_i^2$. If $\lambda_{d_\TS}>\lambda_{d_\TS+1}$, the value under full feature matching is
$\mathcal N+\frac12\sum_{i>d_\TS}\|F^{1/2}\mathsf Cu_i\|^2/\lambda_i$.
\item Suppose $\sigma_{d_\TS}>\sigma_{d_\TS+1}$, and let $t_\sharp:=(v_1,\dots,v_{d_\TS})^\top\Sigma^{-1/2}h_\TT$. For every $\gamma>0$, the
minimisers of $\cL+\gamma\min_A\E\|Ah_\TS-t_\sharp\|^2$ are exactly the output-optimal students.
\item If $h_\TT\sim\cN(0,\Sigma)$ and $g$ is weakly differentiable with $\E\|J\|<\infty$, then $\mathsf C=\bar J\Sigma$ with $\bar J:=\E J(h_\TT)$,
by Stein's identity. So $G_\sharp=F^{1/2}\bar J\Sigma^{1/2}$, and $t_\sharp$ is the target of \cref{def:fisher-rep} with $K$ replaced by
$\bar K:=\bar J^\top F\bar J$.
\item With $K=\E[J^\top FJ]$ itself, the target can cost arbitrarily much. By Jensen, $K\succeq\bar K$, and
$K-\bar K=\E[(J-\bar J)^\top F(J-\bar J)]$ is the Jacobian's variance. Example: $d_\TT=2$, $d_\TS=1$, $h_\TT\sim\cN(0,I)$ and
$g(h)=w_1(h_1^2-1)+w_2h_2$ with $w_2^\top Fw_2<4\,w_1^\top Fw_1$. Then $K=\diag(4w_1^\top Fw_1,\,w_2^\top Fw_2)$ selects $h_1$, which is
uncorrelated with the teacher, at cost $\frac12w_2^\top Fw_2$. $\bar K$ selects $h_2$ at no cost.
\end{enumerate}
For a linear teacher, $\mathsf C=W\Sigma$, $\mathcal N=0$, $G_\sharp=G$ and $K=\bar K$. This recovers \cref{thm:overconstraint,prop:fisher-rep}
and shows that they need no Gaussian assumption.
\end{theorem}
\begin{proof}
With the optimal bias, write $g(h_\TT)-\E g=\mathsf C\Sigma^{-1}h_\TT+r$ with $\E[rh_\TT^\top]=0$. For $L=W_\TS B$ the cross term vanishes, and
$\cL=\mathcal N+\frac12\|F^{1/2}L\Sigma^{1/2}-G_\sharp\|_F^2$. This is the problem of \cref{thm:overconstraint} with $G_\sharp$ in place of $G$.
(a) and (b) follow as in the proofs of \cref{thm:overconstraint}(a,b) and \cref{prop:fisher-rep}, with
$G_\sharp u_i=F^{1/2}\mathsf Cu_i/\sqrt{\lambda_i}$. For (b), every output-optimal $L$ has $F^{1/2}L\Sigma^{1/2}=G_\sharp V_{d_\TS}V_{d_\TS}^\top$
(Eckart--Young, with the gap). This matrix has rank $d_\TS$ and its row space lies in that of $B\Sigma^{1/2}$, so
$B\Sigma^{1/2}=MV_{d_\TS}^\top$ with $M$ invertible, and $h_\TS=Mt_\sharp$.
(c) Stein's identity $\E[g(h)h^\top]=\E[Dg(h)]\,\Sigma$ for $h\sim\cN(0,\Sigma)$ \citep{stein1981estimation}.
(d) $J=(2h_1w_1,\;w_2)$, so $\bar J=(0,w_2)$, and $\E h_1=0$ removes the off-diagonal entries of $K$. Keeping $h_1$:
$\Cov(g,h_1)=\E[(h_1^2-1)h_1]\,w_1=0$, so the best linear student is constant, and
$\cL=\frac12\E\|F^{1/2}g\|^2=w_1^\top Fw_1+\frac12w_2^\top Fw_2$ because $\Var(h_1^2)=2$. Keeping $h_2$: $\cL=\mathcal N=w_1^\top Fw_1$.
\end{proof}

\paragraph{Lossy regime II: nonlinear read-outs.}
With a universal student head, the best head is a conditional expectation. In whitened coordinates
$\xi:=\Sigma^{-1/2}h_\TT$, a target subspace $\mathcal S\subseteq\R^{d_\TT}$ of dimension $\le d_\TS$ therefore costs
\[
\cL^\star(\mathcal S):=\tfrac12\E\big\|F^{1/2}\big(\tilde z_\TT-\E[\tilde z_\TT\mid P_{\mathcal S}\xi]\big)\big\|^2 .
\]
The output optimum solves a nonlinear sufficient-dimension-reduction problem, $\min_{\mathcal S}\cL^\star(\mathcal S)$. Neither
$K$ nor $\bar K$ solves it.

\begin{proposition}[No second-moment Jacobian rule is cost-free]\label{prop:no-moment-rule}
\Mo\Ck Let $d_\TT=2$, $d_\TS=1$ and $h_\TT\sim\cN(0,I)$. For $j\in\{1,2\}$ let $\tilde z^{(j)}_\TT=w_1\psi_j(h_1)+w_2h_2$ with
$w_1^\top Fw_1=1$ and $w_2^\top Fw_2=\frac12$, where $\psi_1(x)=(x^2-1)/\sqrt2$, $\psi_2(x)=c\,(\cos3x-e^{-9/2})$ and
$c^2=4/(9(1-e^{-18}))$. Both teachers have $\Sigma=I$, $\bar J=(0,w_2)$ and $K=\diag(2,\frac12)$. For teacher~1 the unique optimal
target direction is $e_1$, with cost $\frac14$ against $\frac12$ for $e_2$. For teacher~2 it is $e_2$, with cost
$\frac19\frac{1-e^{-9}}{1+e^{-9}}\approx0.111$ against $\frac14$ for $e_1$. Hence every rule that picks the target from
$(\Sigma,\bar J,K)$ has a positive over-constraint cost on one of the two teachers. The $K$-rule loses $\approx0.139$ on
teacher~2, and the $\bar K$-rule loses $\frac14$ on teacher~1. Rescaling $(w_1,w_2)$ by $\lambda$ rescales every cost by $\lambda^2$.
\end{proposition}
\begin{proof}
Each $\psi_j$ is even, with $\E\psi_j=0$ and $\E\psi_j'^2=2$, and only these moments enter $(\Sigma,\bar J,K)$. The cross terms
vanish because $\E\psi_j'(h_1)=0$. For the direction $(\cos\alpha,\sin\alpha)$, let $t=h_1\cos\alpha+h_2\sin\alpha$ and $\rho=\cos\alpha$. Use
$\E[He_k(h_1)\mid t]=\rho^kHe_k(t)$ and $\cos(\omega x)=e^{-\omega^2/2}\sum_{k\ \mathrm{even}}(-1)^{k/2}\omega^kHe_k(x)/k!$. The explained part
$\frac12\E\|F^{1/2}\E[\tilde z_\TT\mid t]\|^2$ equals $\frac12\big[\frac12(1-\rho^2)+\rho^4\big]$ for teacher~1 and
$\frac12\big[\frac12(1-\rho^2)+c^2e^{-9}(\cosh(9\rho^2)-1)\big]$ for teacher~2; the odd--even cross term vanishes. Both are convex in
$\rho^2\in[0,1]$, so each is maximised only at an endpoint. Compare the endpoints.
\end{proof}
\noindent
The mechanism is the Gaussian Poincaré inequality. A degree-$k$ Hermite component of $g$ contributes $k$ times its output
variance to $\tr K$ in whitened coordinates, so $K$ over-weights high-frequency, low-variance nonlinearities; $\bar K$
ignores every nonlinear component. The two still bracket the truth.

\begin{theorem}[Fisher alignment is nearly free: a Poincaré sandwich]\label{thm:poincare}
\Mo\Ck Let $h_\TT\sim\cN(0,\Sigma)$ with $\Sigma\succ0$, let $g$ be weakly differentiable with $\E\|J\|^2<\infty$, let the loss be quadratic
with constant $F$, and let the student head be universal. Put $K_w:=\Sigma^{1/2}K\Sigma^{1/2}$ and $\bar K_w:=\Sigma^{1/2}\bar K\Sigma^{1/2}$, with
$K$ and $\bar K$ built from this $F$, and write $\mu_1(\cdot)\ge\mu_2(\cdot)\ge\cdots$ for eigenvalues. For every subspace $\mathcal S$,
\[
\cL^\star(\mathcal S)=\tfrac12\big\|F^{1/2}\bar J\Sigma^{1/2}P_{\mathcal S^\perp}\big\|_F^2+\mathcal N(\mathcal S),\qquad
0\le\mathcal N(\mathcal S)\le\tfrac12\tr\big(P_{\mathcal S^\perp}(K_w-\bar K_w)\big).
\]
So the linear part of the cost is exact and the nonlinear part is bounded by the Jacobian's variance. If $\mathcal S_K$ is
spanned by the top-$d_\TS$ eigenvectors of $K_w$, which is the target of \cref{def:fisher-rep}, then for every $\gamma>0$
\[
\Delta(\gamma)\le\Delta_\infty\le\tfrac12\sum_{i>d_\TS}\big(\mu_i(K_w)-\mu_i(\bar K_w)\big).
\]
The bound vanishes for linear teachers (\cref{prop:fisher-rep}) and when $k_F\le d_\TS$ (\cref{thm:fisher-central}). Full feature
matching obeys no bound of this kind (\cref{thm:nl-overconstraint}(c)).
\end{theorem}
\begin{proof}
Let $f(\xi):=F^{1/2}g(\Sigma^{1/2}\xi)$ with $\xi\sim\cN(0,I)$. Then $\E[Df^\top Df]=K_w$, and $\E Df=F^{1/2}\bar J\Sigma^{1/2}=:\bar D$
by Stein's identity, so $\bar D^\top\bar D=\bar K_w$. Write $f-\E f=\bar D\xi+r$, where $r$ has no Hermite component of degree
$\le1$ and $\E Dr=0$. Split $\xi=(\xi_\parallel,\xi_\perp)$ into coordinates along $\mathcal S$ and $\mathcal S^\perp$; they are independent.
Conditioning on $\xi_\parallel$ keeps exactly the Hermite terms that do not involve $\xi_\perp$, so
$f-\E[f\mid\xi_\parallel]=\bar DP_{\mathcal S^\perp}\xi+(r-\E[r\mid\xi_\parallel])$, and the two summands are orthogonal because their
Hermite degrees differ. This gives the identity with $\mathcal N(\mathcal S)=\frac12\E\|r-\E[r\mid\xi_\parallel]\|^2$. For fixed
$\xi_\parallel$, the Gaussian Poincaré inequality in $\xi_\perp$ gives
$\E[\|r-\E[r\mid\xi_\parallel]\|^2\mid\xi_\parallel]\le\E[\|Dr\,P_{\mathcal S^\perp}\|_F^2\mid\xi_\parallel]$. Average, and use
$\E[Dr^\top Dr]=K_w-\bar K_w$. For the corollary, the upper bound at $\mathcal S_K$ is $\frac12\tr(P_{\mathcal S_K^\perp}K_w)=\frac12\sum_{i>d_\TS}\mu_i(K_w)$.
By Ky Fan, the lower bound for any $\mathcal S$ of dimension $\le d_\TS$ is at least $\frac12\sum_{i>d_\TS}\mu_i(\bar K_w)$. Finish with
\cref{lem:penalty}(iii).
\end{proof}
\noindent
The upper bound is the gradient-based dimension-reduction bound of \citet{zahm2020gradient}, applied to $F^{1/2}g$. The
lower bound is the linear (Stein) part. The gap between them is the only obstruction to a cost-free spectral target,
and \cref{prop:no-moment-rule} shows that the gap can be crossed in either direction.

% ---------------------------------------------------------------------------
\subsection{Sample complexity and computation of intermediate supervision}\label{sec:repr-complexity}
% ---------------------------------------------------------------------------

Let the teacher compute through stages $c_\ell=\phi_\ell(s,c_{<\ell})$, $\ell=1,\dots,L$, with output $y=c_L$. A stage can be a
layer, in which case $c_\ell$ is a representation that depends on $c_{\ell-1}$ only, or a chain-of-thought step, in which
case $c_\ell$ is an emitted token or segment. \emph{Intermediate supervision} (IS) reveals $c_{1:L}(s_i)$ on the training
states; \emph{end-to-end} (E2E) training reveals only $y(s_i)$. The student learns $\hat\phi_\ell\in\cH_\ell$ and computes
$\hat c_\ell=\hat\phi_\ell(s,\hat c_{<\ell})$.

\begin{theorem}[Composition]\label{thm:composition}
\Pa{realisability, $\phi_\ell\in\cH_\ell$} Let $e_\ell:=\P_{s\sim\cD}[\hat\phi_\ell(s,c_{<\ell}(s))\neq c_\ell(s)]$ be the error of stage $\ell$ on
\emph{teacher} prefixes.
\begin{enumerate}[label=(\roman*),nosep]
\item $\P_{s\sim\cD}[\hat y(s)\neq y(s)]\le\sum_{\ell=1}^Le_\ell$. Errors add; they do not compound.
\item Let the $\cH_\ell$ be finite and let each $\hat\phi_\ell$ be consistent with the $n$ supervised examples of its stage. With
probability $\ge1-\delta$, $\P[\hat y\neq y]\le\frac1n\sum_\ell\log\frac{L|\cH_\ell|}{\delta}$. The running time is $\sum_\ell T_\ell(n)$,
where $T_\ell$ is the time of a consistent learner for $\cH_\ell$.
\item E2E, any hypothesis in the composed class $\cH_\circ$ that is consistent with the outputs satisfies
$\P[\hat y\neq y]\le\frac1n\big(\sum_\ell\log|\cH_\ell|+\log\frac1\delta\big)$. Finding one can be intractable even when every $\cH_\ell$
has a polynomial-time consistent learner (\cref{prop:parity,prop:graded}(b)).
\end{enumerate}
\end{theorem}
\begin{proof}
(i) If every stage is correct on teacher prefixes at $s$, induction on $\ell$ gives $\hat c_\ell(s)=c_\ell(s)$ for all $\ell$. This
is the first-deviation argument of \cref{prop:hybrid} for a deterministic teacher (\cref{ex:regimes}(b)).
(ii) A fixed $\phi\in\cH_\ell$ with error $>\eps_\ell$ is consistent with $n$ examples with probability $<(1-\eps_\ell)^n\le e^{-n\eps_\ell}$.
Take a union bound over $\cH_\ell$ and over $\ell$, with $\eps_\ell=\log(L|\cH_\ell|/\delta)/n$, and use (i).
(iii) The same argument for the single class $\cH_\circ$, with $|\cH_\circ|\le\prod_\ell|\cH_\ell|$.
\end{proof}

\noindent
The capacity bounds in (ii) and (iii) have the same order. In the realisable, noiseless setting, capacity arguments
therefore cannot show that IS helps statistically. Its gains come from four other sources: noise in the output
channel (\cref{prop:repr-count}), computation (iii), restricted learners such as gradient methods
(\cref{prop:graded}(a)), and the information that each intermediate value carries per sample. The last can be
unbounded:

\begin{example}[An unbounded statistical gain without noise]\label{ex:retrieval}
\Pv\Ck Let $\cD$ be uniform on $[N]$ with $N\ge8$. The teacher first retrieves a key, $c_1=k\in[N]$, which is a constant function
of $s$, and then outputs $y=\ind\{s=c_1\}$. Let $\cH_1$ be the constant maps and let $\cH_2=\{\phi_2\}$. With IS, one example reveals
$k$, and the student is exact. E2E, any hypothesis with error below $\frac1N$ must equal $\ind\{\cdot=k\}$. If the sample misses
$k$, which happens with probability $(1-\frac1N)^n$ when $k$ is uniform, the learner has seen only zeros and can do no better
than guess among at least $N-n$ unseen values. So every E2E learner with $n\le N/2$ examples fails to reach error below $\frac1N$
with probability at least $(1-\frac1N)^{N/2}(1-\frac2N)\ge\frac14$ for some $k$. The intermediate value reveals in one example what
the output reveals only on an event of probability $\frac1N$.
\end{example}

\paragraph{Transformer layers.}
For attention the supervised quantities are the attention maps, as in MiniLM and TinyBERT
\citep{wang2020minilm,jiao2020tinybert}, and the residual stream.

\begin{proposition}[Attention-map supervision makes layerwise distillation convex]\label{prop:attention-convex}
\Pv\Ck Consider one attention head on a context $X=(x_1,\dots,x_N)\in\R^{d\times N}$ with query token $x_q$. Its attention is
$\alpha_Q(\cdot\mid X)=\softmax_j(x_q^\top Qx_j)$ and its output is $o=VX\alpha_Q$, with $Q=W_Q^\top W_K\in\R^{d\times d}$ and $V\in\R^{d'\times d}$
unconstrained (the merged parametrisation). The teacher uses $(Q^\circ,V^\circ)$.
\begin{enumerate}[label=(\roman*),nosep]
\item Given the teacher's attention maps on contexts $X_1,\dots,X_n$, the objective
$\sum_c\KL\big(\alpha_{Q^\circ}(\cdot\mid X_c)\,\|\,\alpha_Q(\cdot\mid X_c)\big)$ is convex in $Q$. Given the attention maps, fitting the
outputs is least squares in $V$.
\item If the matrices $x_{q,c}(x_{j,c}-x_{j',c})^\top$ span $\R^{d\times d}$, every minimiser of the first objective reproduces the
teacher's attention on \emph{every} context. For contexts drawn from a law with a density this holds almost surely once
$n\ge d\,\lceil d/\min(N-1,d)\rceil$. If the vectors $X_c\alpha_c$ span $\R^d$, which holds almost surely once $n\ge d$, then $V=V^\circ$.
\item The end-to-end objective $\sum_c\|VX_c\alpha_Q(\cdot\mid X_c)-o_c\|^2$ is not convex in $(Q,V)$: its zero set is not convex.
\end{enumerate}
\end{proposition}
\begin{proof}
(i) The attention logits $\langle Q,x_qx_j^\top\rangle$ are linear in $Q$, and $z\mapsto\KL(p\|\softmax z)$ is convex
(\cref{prop:convexity}(i)). (ii) The minimum value is $0$, attained at $Q^\circ$. A zero means that the logits agree up to a
per-context constant, that is, $\langle Q-Q^\circ,x_{q,c}(x_{j,c}-x_{j',c})^\top\rangle=0$ for all $c,j,j'$, and spanning forces $Q=Q^\circ$.
The Gram determinant of the spanning family is a polynomial in the context entries, so it vanishes on a null set unless
it vanishes identically. It does not: let the queries run through $e_1,\dots,e_d$, each repeated $\lceil d/m\rceil$ times with
$m=\min(N-1,d)$, and let the key differences run through a basis of $\R^d$ in blocks of $m$. The map
$(X_c)_c\mapsto\det(X_1\alpha_1,\dots,X_d\alpha_d)$ is real-analytic and equals $1$ when every token of $X_c$ is $e_c$, so it too is
nonzero almost surely. (iii) Take $d=d'=1$, $N=2$, $x_q=1$, one context $(1,2)$ and $o\neq0$. The zero set is the curve
$V=o/m(Q)$ with $m(Q)=2-\sigma(-Q)$, which is strictly monotone and not affine, so the set is not convex.
\end{proof}
\noindent
With every layer's attention maps, residual stream and MLP pre-activations supervised, a student with the teacher's
architecture is therefore distilled exactly by $L$ rounds of convex programs; the rank constraint $d_h<d$ of a smaller
head dimension is what breaks convexity. \Cref{thm:composition}(i) is the discrete analogue of the fact that exact
stages compose exactly. To quantify gains \emph{beyond parities} we need end-to-end lower bounds that grade with the
target, or that hold for all efficient algorithms.

\begin{proposition}[Separations beyond parities]\label{prop:graded}
\Kn
\begin{enumerate}[label=(\alph*),nosep]
\item (\emph{Gradient methods: a polynomial hierarchy.}) Let $x\sim\cN(0,I_d)$. Let the teacher compute the pre-activation
$h_\TT=\langle w,x\rangle$, for example one MLP unit or one attention logit, and output $g(h_\TT)$, where the link $g$ has
information exponent $k^\star$: the index of its first non-zero Hermite coefficient. From outputs alone, correlational
statistical-query algorithms, which include gradient descent on the square loss, need query tolerance
$\tau\lesssim d^{-k^\star/4}$, that is $n\gtrsim d^{k^\star/2}$ samples, and online SGD needs $\tilde\Theta(d^{\max(k^\star-1,1)})$ samples
\citep{benarous2021online,damian2023smoothing}. With $h_\TT$ as an intermediate target, least squares recovers $w$ exactly from
$n=d$ noiseless samples in general position, or to squared error $\approx\sigma^2d/n$ from targets with noise variance $\sigma^2$,
and what remains is a one-dimensional problem. The gain factor is $d^{\max(k^\star/2-1,0)}$: nothing for $k^\star\le2$, and a
polynomial that grows with $k^\star$ beyond that. Multi-index teachers behave alike, with the leap complexity in place of
$k^\star$ \citep{abbe2023sgd}. Sparse parities are the Boolean analogue, with the degree in place of $k^\star$.
\item (\emph{All efficient algorithms.}) Under standard lattice assumptions, no polynomial-time algorithm learns intersections of
$d^{\epsilon}$ halfspaces from labelled examples \citep{klivans2009cryptographic}. Such a function is a two-layer threshold network,
hence a single MLP block. With its hidden units supervised, each unit is a realisable halfspace, learnable in polynomial
time by linear programming, and \cref{thm:composition}(i) bounds the end-to-end error by the sum of the unit errors.
\end{enumerate}
\end{proposition}
\noindent
Intermediate supervision by chain-of-thought tokens also changes the approximation term. Constant-depth,
log-precision transformers compute only functions in uniform $\mathsf{TC}^0$ \citep{merrill2023parallelism}, while with
polynomially many intermediate tokens they can simulate polynomial-time computation
\citep{feng2023towards,merrill2024expressive,li2024chain}. Distilling a teacher's intermediate tokens can therefore lower
$\cA$ from $\Omega(1)$ to $0$ for a fixed-depth student. Learning to emit them is then a realisable sequence-level problem, with
the horizon-free rate of \cref{thm:mle}; \citet{malach2024autoregressive} develops this learnability view for next-token
predictors trained on such data.
